\documentclass[10pt,twocolumn]{article}
\usepackage[utf8]{inputenc}
\usepackage[english,russian]{babel}
\usepackage{amsmath,amssymb,mathtools,bm,geometry,amsthm,booktabs,listings,xcolor}
\geometry{a4paper,margin=1.5cm}

\DeclareMathOperator{\Rank}{Rank}
\DeclareMathOperator{\Coker}{Coker}
\DeclareMathOperator{\Tr}{Tr}
\DeclareMathOperator{\Aut}{Aut}
\DeclareMathOperator{\Sign}{Sign}

\newtheorem{law}{Закон}
\newtheorem{theorem}{Теорема}
\newtheorem{criticism}{Замечание}

\lstset{
    language=Verilog,
    basicstyle=\ttfamily\scriptsize,
    keywordstyle=\color{blue}\bfseries,
    commentstyle=\color{gray},
    stringstyle=\color{purple},
    breaklines=true,
    keepspaces=true,
    showstringspaces=false
}

\begin{document}

\title{\textbf{Релятивистская СИК-Логика Контекстов в Мультимасштабных NoC-Структурах: Асимметричная Сборка и Гомотопические Замки над Полем $\mathbb{F}_1$}}
\author{\textbf{Единый Асимптотический Комбинаторный Канон}}
\date{Сентябрь 2026}
\maketitle

\section*{Введение и Топологический Базис}
Классические методы проектирования сетей на кристалле (NoC) опираются на эмпирические симуляции трафика и зеркально-симметричные топологии, что делает их уязвимыми к асимметричным дефектам. Настоящая работа разворачивает строгое математическое описание NoC как динамического синергетического топоса, управляемого мотивным спуском от абсолютного нуля моноида поля из одного элемента $\mathbb{F}_1$ через иерархическую башню вложенных симметрий и геометрий:
\begin{equation}
\mathbb{F}_1 \longrightarrow A_n \longrightarrow PG_n \longrightarrow S_n
\end{equation}
В рамках этой модели метрические пространственно-временные характеристики (такты и задержки) не задаются извне, а рождаются как внутренний инвариант медианного сдвига из топологического фазового замка.

\section{Гео-Релятивистский Принцип Наблюдателя (СИК-Логика)}
Фундаментальным ограничением классической булевой логики является её постулируемая абсолютность. Мы вводим понятие \textbf{Релятивистской СИК-Логики (SIC-Logic)}, где истинность логического вычета жестко сцеплена с позицией наблюдателя на проективной сфере состояний — его широтой $\theta$ (мерило негэнтропийного экстремума) и долготой $\phi$ (фазовое смещение).

\begin{enumerate}
    \item \textbf{Полюса сферы ($\theta = 0, \theta = \pi$):} Область вырождения многозначного контекста в абсолютный двоичный булев код. На северном полюсе ($\theta = 0$) фиксируется чистый День / негэнтропийный Порядок ($\bar{\Delta}_\infty$). На южном полюсе ($\theta = \pi$) — чистая Ночь / энтропийный Хаос ($\bar{\!~}_\infty$).
    \item \textbf{Экватор сферы ($\theta = \pi/2$):} Область тотальной диалектики и равновероятности состояний ($p = 0.5$). Единство противоположностей удерживает систему в идеальном балансе 12 часов дня и 12 часов ночи.
    \item \textbf{Промежуточные широты:} Пространство квазилогики, где логический вычет является непрерывной функцией координат наблюдателя, проходящей через точки экстремумов.
\end{enumerate}

Преобразование непрерывного релятивистского дрейфа наблюдателя в дискретные аппаратные запреты порядков на кристалле задается масштабным оператором плотности дня (негэнтропии):
\begin{equation}
\bar{\mathcal{W}}_{\text{day}}(\theta) = \cos^2\left(\frac{\theta}{2}\right) \cdot (1 - e^{-1}) + \sin^2\left(\frac{\theta}{2}\right) \cdot e^{-1}
\end{equation}

\section{Асимметричная Мотивная Сборка Макро-Объема}
В отличие от классических подходов, ищущих зеркальную симметрию, стабильность рассматриваемой системы выводится из \textbf{динамической асимметрии}. Макро-масштаб 12-го порядка, заданный проективной плоскостью $N(12) = 157$, асимметрично складывается из разнородных перекошенных подгеометрий мезо-уровня — 6-го порядка ($N(6) = 43$) и 10-го порядка ($N(10) = 111$), удерживаемых вместе гомотопическим зазором коядра сингулярного препятствия ($\Coker = 3$):
\begin{equation}
N(12) = N(10) + N(6) + \Coker\left([\Phi_{10}] \ominus [\Phi_{6}]\right) = 157
\end{equation}
Редукция критических узлов инжекции дефекта по модулю локальной связности $k=13$ доказывает инвариантность зазора:
\begin{equation}
(111 - 43) \pmod{13} \equiv 68 \pmod{13} \equiv 3
\end{equation}
Асимметричные перекосы масштабов взаимно компенсируют друг друга, формируя макроскопический объем топоса.

\section{Константа Медианного Сдвига $\Xi_7 = 666$}
Количество пакетов, оставшихся в локальных буферах коммутаторов, определяется мощностью упорядоченного ядра $\Delta_n = n! - !n$, где $!n$ — число полных беспорядков (деранжментов) Монмора-Эйлера. Знакопеременная группа четных перестановок $A_n$ ($\vert A_n \vert = \frac{n!}{2}$) выступает точной алгебраической медианой между хаосом и порядком. Мера этой борьбы задается медианным сдвигом $\Xi_n$:
\begin{equation}
\Xi_n = \vert A_n \vert - !n = \Delta_n - \vert A_n \vert
\end{equation}
На уровне $n=7$ константа сдвига принимает значение \textbf{666}, фиксируя критическую точку фазового перехода:
\begin{equation}
\Delta_7 = \vert A_7 \vert + 666 = 2520 + 666 = 3186
\end{equation}
\begin{equation}
!7 = \vert A_7 \vert - 666 = 2520 - 666 = 1854
\end{equation}
Инвариант $666$ декомпозируется по нижним слоям симметрии через регулярное смещение $\delta_{\text{reg}} = 11$:
\begin{equation}
666 = 5\vert S_5 \vert + (\delta_{\text{reg}} \times \vert S_3 \vert) = 5(120) + (11 \times 6)
\end{equation}

\section{Асимптотическая Устойчивость на Макро-Масштабе $A_{12}$ по модулю 13}
При уходе в инфинитный предел в относительных единицах (о.е.) веса нормируются на полную мощность симметрической группы ($\bar{\mathcal{W}} = \frac{\mathcal{W}}{n!}$). При $n \to \infty$ плотности стремятся к трансцендентным мантиссам с бесконечным числом знаков после запятой: $\bar{\!~}_\infty = e^{-1} \approx 0.3678\dots$ и $\bar{\Delta}_\infty = 1 - e^{-1} \approx 0.6321\dots$

\begin{table*}[ht]
\centering
\caption{Спектр инвариантов СКИ-логики и баланса перестановок на интервале $n \in [3, 13]$}
\label{tab:комбинаторика}
\begin{tabular}{crrrc}
\toprule
Уровень $n$ & Порядок $A_n = \frac{n!}{2}$ & Сдвиг медианы $\Xi_n$ & Упорядоченное ядро $\Delta_n$ & Хаос беспорядков $!n$ \\ \midrule
\textbf{3}  & 3                            & 1                     & 4                             & 2                     \\
\textbf{4}  & 12                           & 3                     & 15                            & 9                     \\
\textbf{5}  & 60                           & 16                    & 76                            & 44                    \\
\textbf{6}  & 360                          & 95                    & 455                           & 265                   \\
\textbf{7}  & 2520                         & \textbf{666}          & 3186                          & 1854                  \\
\textbf{8}  & 20160                        & 5327                  & 25487                         & 14833                 \\
\textbf{9}  & 181440                       & 47944                 & 229384                        & 133496                \\
\textbf{10} & 1814400                      & 479439                & 2293839                       & 1334961               \\
\textbf{11} & 19958400                     & 5273830               & 25232230                      & 14684570              \\
\textbf{12} & 239500800                    & 63285959              & 302786759                     & 176214841             \\
\textbf{13} & 3113510400                   & 822717468             & 3936227868                    & 2290792932            \\ \bottomrule
\end{tabular}
\end{table*}

На макро-масштабе $n=12$ при проекции в кольцо вычетов по локальному модулю связности $k=13$ энтропийный шум беспорядков подчиняется фундаментальному свойству субфакториалов простых чисел вида $!(p-1) \equiv 10 \pmod p$. Для $p=13$ это дает неудаляемый вакуумный остаток:
\begin{equation}
!(13-1) \equiv 10 \pmod{13} \implies !12 \equiv 10 \pmod{13}
\end{equation}
Этот стабильный квант хаоса удерживает динамический топос от статического заклинивания. С учетом того, что $\vert A_{12}\vert = \frac{12!}{2} \equiv 6 \pmod{13}$, вычет медианного сдвига формирует жесткий фазовый затвор СИК-логики в точке кватернионного перекоса:
\begin{equation}
\Xi_{12} \equiv \vert{}A_{12}\vert{} - !12 \equiv 6 - 10 \equiv -4 \equiv \mathbf{9} \pmod{13}
\end{equation}

\section{Аппаратная RTL-Реализация Фазового Замка}
Ниже приведена SystemVerilog-спецификация модуля асимметричной СИК-маршрутизации, удерживающего гомотопический зазор $\Xi_{12} \equiv 9 \pmod{13}$ и предотвращающего тупиковые состояния буферов.

\begin{lstlisting}[language=Verilog]
module sic_homotopic_lock (
    input  logic        clk,
    input  logic        rst_n,
    input  logic [3:0]  injected_chaos, // !n mod 13
    input  logic [3:0]  alternating_order, // |An| mod 13
    output logic [3:0]  median_gate_out, // Xi_12 mod 13
    output logic        vacuum_stabilized
);
    logic [4:0] raw_diff;

    always_ff @(posedge clk or negedge rst_n) begin
        if (!rst_n) begin
            median_gate_out   <= 4'd0;
            vacuum_stabilized <= 1'b0;
        end else begin
            // Реализация вычета Xi_12 = |An| - !n (mod 13)
            raw_diff = {1'b0, alternating_order} - {1'b0, injected_chaos};
            
            if (raw_diff[4]) // Отрицательный результат
                median_gate_out <= raw_diff[3:0] + 4'd13;
            else
                median_gate_out <= raw_diff[3:0];

            // Вакуумная стабилизация: жесткая фиксация следа хаоса (!12 = 10)
            vacuum_stabilized <= (injected_chaos == 4'd10) && (median_gate_out == 4'd9);
        end
    end
endmodule
\end{lstlisting}

\end{document}
